% fig_bracketing.tex -- NEW native pgfplots (figure*). The mechanism of Theorem 1:
% finitely many moments BRACKET the cumulative weight F(t); the bracket width IS W_n(t), contracting with n.
% 2026-09-27 layout fix: (i) clip mode=individual, so the t* label sits just above the top frame instead of being
% clipped under it; (ii) the inset is typeset into a box first and placed on the main axis's FOREGROUND layer on an
% opaque white panel inside the empty lower-left region of the main axis (t in [-0.955,-0.70], F in [0.14,0.66]:
% F(t) < 0.03 and both brackets < 0.1 there), so the main grid (drawn on top by axis on top + fillbetween layers)
% no longer shows through the inset and its y label no longer sits on the main 0.4 tick label;
% (iii) inset ymax 1.04 so the n=1 marker (1.000) clears the inset's top frame; (iv) main grid drawn below the data.
% ---- inset: max bracket width vs order (same christoffel_maxW.dat as fig_christoffel (b)) ----
\ifdefined\bracketinsetbox\else\newsavebox\bracketinsetbox\fi
\sbox\bracketinsetbox{%
\begin{tikzpicture}
\begin{axis}[
  width=4.1cm, height=2.95cm,
  paperaxis, grid=none,
  xmin=0.7, xmax=4.3, ymin=0.82, ymax=1.04, xtick={1,2,3,4}, ytick={0.85,0.95},
  xlabel={order $n$}, ylabel={$\max_t W_n$},
  xlabel style={font=\scriptsize}, ylabel style={font=\scriptsize},
  tick label style={font=\scriptsize},
]
  \addplot[pOx, very thick, mark=*, mark size=1.8pt, mark options={fill=pAmber,draw=pOx}]
      table[x=n,y=maxW]{figs/christoffel_maxW.dat};
\end{axis}
\end{tikzpicture}}%
\begin{tikzpicture}
\begin{axis}[
  name=main, paperaxis,
  width=12.4cm, height=6.9cm,
  xmin=-1, xmax=-0.35, ymin=0, ymax=1.03,
  xtick={-1,-0.9,-0.8,-0.7,-0.6,-0.5,-0.4},
  % F(t) and the brackets run to t=1; each plot is clipped to the axis box (clip mode=individual), while the
  % annotation nodes below are not, so the t* label can sit above the top frame
  xlabel={rescaled frequency $t$}, ylabel={cumulative weight $F(t)=\int^{t}\!A$},
  clip=true, clip mode=individual,
  axis on top=false,   % grid below the data: on top, the F=1 grid line split the thick staircase into two thin lines
]
  % n=2 bracket (light amber) = extremal spread over measures sharing m_0..m_4
  \addplot[name path=lo2, draw=none, forget plot] table[x=t,y=Flo]{figs/bracket_env_n2.dat};
  \addplot[name path=hi2, draw=none, forget plot] table[x=t,y=Fhi]{figs/bracket_env_n2.dat};
  \addplot[pAmber, fill opacity=0.26, forget plot] fill between[of=lo2 and hi2];
  % n=4 bracket (deep oxblood) nests inside
  \addplot[name path=lo4, draw=none, forget plot] table[x=t,y=Flo]{figs/bracket_env_n4.dat};
  \addplot[name path=hi4, draw=none, forget plot] table[x=t,y=Fhi]{figs/bracket_env_n4.dat};
  \addplot[pOx, fill opacity=0.34, forget plot] fill between[of=lo4 and hi4];
  % the true cumulative weight lies inside every bracket
  \addplot[pInk, very thick, const plot] table[x=t,y=F]{figs/bracket_F.dat};

  % interior support point t* (dominant atom): guide inside, label just above the top frame (outside the data)
  \draw[guide] (axis cs:-0.5795,0) -- (axis cs:-0.5795,1.0);
  \node[pInk, font=\scriptsize, anchor=south, inner sep=1.5pt] at (axis cs:-0.5795,1.03) {$t^{*}$};

  % labels (short, in the opened whitespace)
  \node[pAmber!75!black, font=\small, anchor=east] at (axis cs:-0.62,0.55) {$n{=}2$};
  \node[pOx, font=\small, anchor=east] at (axis cs:-0.62,0.40) {$n{=}4$};
  \node[pInk, font=\small, anchor=east] at (axis cs:-0.38,0.92) {true $F(t)$};
  \node[pInk, font=\scriptsize, anchor=north west, align=left] at (axis cs:-0.99,0.98)
      {finitely many moments do not\\ pin $F(t)$ --- they \emph{bracket} it;\\ the band width is the\\ Christoffel function $W_n(t)$};

  % the inset, on an opaque white panel above the main grid
  \pgfplotsonlayer{axis foreground}
    \node[fill=white, inner sep=1pt, anchor=south west] at (axis cs:-0.955,0.14) {\usebox\bracketinsetbox};
  \endpgfplotsonlayer
\end{axis}
\end{tikzpicture}
